Creative thinking stands as one of the most distinctive and complex capacities of the human mind, underpinning scientific discovery, artistic expression, and everyday problem-solving \citep{torrance1974torrance, guilford1967nature, duncker1948problem}. Within the broader landscape of human cognition, creativity is commonly decomposed into two complementary modes: \textit{convergent thinking}, which involves narrowing down multiple possibilities to arrive at a single, well-defined solution, and \textit{divergent thinking}, which involves generating a wide range of novel ideas or associations from a single starting point \citep{guilford1967nature, mednick1962associative}. While both modes contribute to creative cognition, divergent thinking is widely considered the more foundational driver of creativity, as it captures the generative, open-ended exploration that gives rise to original thought \citep{runco2012divergent}. Its study has therefore attracted greater attention in both cognitive psychology and neuroscience, with researchers seeking to understand not only what makes humans capable of such flexible ideation, but also what neural substrates and cognitive mechanisms support it \citep{beaty2018robust, abraham2013promises, beaty2016creative}.

\begin{figure}[t]
\includegraphics[width=\linewidth,trim={1cm 1cm 0cm 1cm}]{figures/cadabra_method_v2.png}
\centering
\caption{\textbf{Our high-level brain alignment methodology.}}
\label{fig:method}
\end{figure}

In recent years, large language models (LLMs) have rapidly expanded their reach across a wide range of cognitive tasks, with divergent thinking among the abilities they now exhibit \citep{ismayilzada2024creativity, zhao2023survey}. Models have shown impressive results on established creativity benchmarks such as the Alternative Uses Test (AUT), which measures the ability to generate diverse and unconventional uses for everyday objects, in some cases reaching or exceeding average human performance. \citep{goes2023pushing, stevenson2022putting, bellemare2026divergent}. Alongside these behavioral results, a growing body of work has established that LLMs with higher task performance tend to exhibit greater alignment with the human language system \citep{schrimpf2021neural, goldstein2022shared, caucheteux2022brains}. However, prior work has predominantly focused on passive language processing tasks, in which participants read or listen to naturalistic text \citep{aw2023instruction, alkhamissi2025language, lebel2023natural} or, recently, simple word association \citep{kwon2024assessing} and abstract reasoning tasks \citep{pinier2025large}. Divergent thinking, however, is a higher-order cognitive function often employing a broad and distributed network of brain regions well beyond the classical language system \citep{liu2024neural}. Whether more creative LLMs are correspondingly more aligned to the neural representations underlying human creative cognition, therefore, remains an open and compelling question, and one that existing brain-LLM alignment frameworks have yet to address.

In this work, we take a first step toward addressing this question by systematically investigating brain-LLM alignment in the context of creative thinking using the AUT task. We leverage fMRI data collected from \citet{beaty2018robust}, in which participants performed both a creative task, namely the AUT, and a matched non-creative control task, namely the Object Characteristics Task (OCT), which involves generating typical characteristics of everyday objects \citep{beaty2018robust}. Similarly, we extract model representations to the same stimuli in these tasks from a range of LLMs and measure their alignment to the corresponding brain responses using Representational Similarity Analysis (RSA) \citep{kriegeskorte2008representational}. On the neural side, we examine alignment across several canonical regions of interest involved in creative thinking, namely, the Default Mode Network (DMN) and the Frontoparietal Network (FPN) \citep{luchini2025role, liu2024neural}. 
% We further complement this atlas-based approach with functional localization \citep{fedorenko2010new}, a technique borrowed from cognitive neuroscience, to identify a creative model network, and examine alignment between these functionally defined networks and brain regions. 
Our results reveal that brain alignment during creative thinking is stage-dependent: alignment scales positively with model size and creative task performance when measured at the prompt stage, but this relationship weakens once models generate responses. We further find that post-training objectives shape alignment in functionally selective and interpretable ways, with creativity-optimized and reasoning-chain training steering model representations toward fundamentally different regions of the neural geometry of human creative thought. To the best of our knowledge, this is the first study to directly examine brain-LLM alignment in the context of an active creative thinking task, and we hope our findings encourage deeper investigation into the neural plausibility of creative cognition in language models. \textbf{Our main contributions are as follows:}
\begin{itemize}
\item We present the first systematic investigation of brain-LLM alignment in the context of an active creative thinking task, extending the existing alignment literature beyond passive language processing to a higher-order cognitive domain.
\item We show that brain alignment during creative thinking scales positively with model size and task performance when measured at the prompt stage, but that this relationship weakens once models generate responses, revealing an important stage-dependent limitation of current alignment analyses.
\item We show that post-training objectives shape alignment in a functionally selective way: creativity-optimized training selectively increases alignment with high-creativity neural responses while decreasing alignment with low-creativity ones, whereas reasoning-chain training produces a striking reversal, suggesting that different training objectives steer model representations toward fundamentally different regions of the neural geometry of human creative thought.
\end{itemize}
